\begin{helper}
Let \(N\) be a natural number, and let
\(b,y,y^\downarrow\in\mathbb R^{\operatorname{Fin}(N)}\).  Suppose that
\(b\) is nonincreasing, that \(y^\downarrow\) is nonincreasing, and that
\(y^\downarrow\) is a permutation of \(y\).  Then sorting \(y\) into the
nonincreasing order \(y^\downarrow\) does not increase its distance from
\(b\):
\[
\sum_k |y^\downarrow_k-b_k|\le \sum_k |y_k-b_k|.
\]
\end{helper}
\begin{proof}
Write
\[
D(u)=\sum_k |u_k-b_k|.
\]
We first prove the two-coordinate inequality used at every adjacent swap.
Let \(a\le d\) and let \(\beta\le\alpha\).  Define
\[
\phi(t)=|d-t|-|a-t|.
\]
The function \(\phi\) is nonincreasing.  Indeed, if \(t\le a\), then
\[
\phi(t)=(d-t)-(a-t)=d-a.
\]
If \(a\le t\le d\), then
\[
\phi(t)=(d-t)-(t-a)=d+a-2t,
\]
which is nonincreasing in \(t\).  If \(d\le t\), then
\[
\phi(t)=(t-d)-(t-a)=a-d.
\]
The three formulas agree at the interval endpoints, and the middle formula
decreases from \(d-a\) to \(a-d\); hence the whole function is
nonincreasing.  Since \(\beta\le\alpha\), we get
\(\phi(\alpha)\le\phi(\beta)\), namely
\[
|d-\alpha|-|a-\alpha|
\le
|d-\beta|-|a-\beta|.
\]
Rearranging gives
\[
|d-\alpha|+|a-\beta|\le |a-\alpha|+|d-\beta|.
\]

Now suppose a vector \(u\) has an adjacent inversion at positions \(r,r+1\):
\[
u_r=a<d=u_{r+1}.
\]
Because \(b\) is nonincreasing, \(b_r\ge b_{r+1}\).  Apply the preceding
inequality with \(\alpha=b_r\) and \(\beta=b_{r+1}\).  If \(u'\) is obtained
from \(u\) by swapping the two adjacent entries \(a,d\), then all terms in
\(D(u')\) and \(D(u)\) are identical except at \(r,r+1\), and the
two-coordinate inequality gives
\[
D(u')\le D(u).
\]
Thus moving a larger adjacent entry leftward never increases the distance to
the sorted vector \(b\).

Starting from \(u^{(0)}=y\), repeatedly swap an adjacent inversion whenever
one exists.  Each step is a transposition of two coordinates, so after any
finite number of steps the current vector is a permutation of \(y\).  The
process terminates because the inversion count
\[
I(u)=|\{(r,s): r<s\text{ and }u_r<u_s\}|
\]
strictly decreases at each adjacent-inversion swap.  To see this, suppose
the swap is at \(p,p+1\), with \(a=u_p<d=u_{p+1}\).  The pair
\((p,p+1)\) contributes \(1\) to \(I(u)\) before the swap and \(0\) after.
For an index \(k<p\), the two possible inversion contributions involving
\(k,p,p+1\) are, before the swap,
\[
\mathbf 1_{u_k<a}+\mathbf 1_{u_k<d},
\]
and, after the swap,
\[
\mathbf 1_{u_k<d}+\mathbf 1_{u_k<a},
\]
the same sum.  For an index \(k>p+1\), the corresponding contributions are
\[
\mathbf 1_{a<u_k}+\mathbf 1_{d<u_k}
\]
before the swap and
\[
\mathbf 1_{d<u_k}+\mathbf 1_{a<u_k}
\]
after the swap, again the same sum.  Pairs not involving \(p\) or \(p+1\)
are unchanged.  Hence \(I\) decreases by exactly \(1\).  Since \(I(u)\) is a
nonnegative integer, only finitely many swaps can occur.

Let \(u^*\) be the terminal vector.  It has no adjacent inversion, so it is
nonincreasing: if every adjacent pair satisfies \(u^*_r\ge u^*_{r+1}\), then
for \(r<s\) the chain
\[
u^*_r\ge u^*_{r+1}\ge\cdots\ge u^*_s
\]
gives \(u^*_r\ge u^*_s\).  From the previous paragraphs, \(u^*\) is a
permutation of \(y\), and
\[
D(u^*)\le D(y).
\]

It remains to compare this terminal sorted permutation with the specified
sorted permutation \(y^\downarrow\).  Both \(u^*\) and \(y^\downarrow\) are
nonincreasing permutations of \(y\), so they have the same multiset of
entries.  We claim they agree coordinate by coordinate.  If not, let \(r\)
be the first position where they differ.  If
\(u^*_r<y^\downarrow_r\), put \(t=y^\downarrow_r\).  Since
\(y^\downarrow\) is nonincreasing, at least the first \(r+1\) entries of
\(y^\downarrow\) are \(\ge t\).  Since \(u^*\) is nonincreasing and
\(u^*_r<t\), at most the first \(r\) entries of \(u^*\) are \(\ge t\).  This
contradicts equality of multisets.  The case \(u^*_r>y^\downarrow_r\) is
the same argument with \(u^*\) and \(y^\downarrow\) interchanged.  Therefore
\(u^*=y^\downarrow\).

Consequently,
\[
\sum_k |y^\downarrow_k-b_k|
=D(y^\downarrow)
=D(u^*)
\le D(y)
=\sum_k |y_k-b_k|,
\]
which is the desired inequality.
\end{proof}
